\bookchapter{פרק 5: תן לטלפון להבין אותך טוב יותר – הגדרות בסיסיות}

\chapterintro{השקעת כמה דקות בהגדרות משפרת את החוויה משמעותית. ההגדרות האלה הן כמו “כוונון” לטלפון שלך, כדי שיתאים יותר להרגלי השימוש שלך.}

\booksection{צליל ורטט1. צליל ורטט}

: בחר צלצול שאתה שומע ושאינו מביך.

\begin{itemize}
  \item \textbf{מצב רטט}: התראות באמצעות רטט כשהטלפון במצב שקט.
  \item \textbf{מצב 'נא לא להפריע'}: ניתן להגדיר חריגים, למשל לאפשר שיחות מאנשי קשר.
  \item \textbf{הפרדת עוצמת שמע}: ניתן לכוון בנפרד עוצמת מדיה, צלצול ושעון מעורר.
  \item \textbf{2. תצוגה ובהירות}בהירות אוטומטית
\end{itemize}

\booksection{תצוגה ובהירות2. תצוגה ובהירות}

\begin{itemize}
  \item \textbf{מצב הגנה על העיניים / מצב לילה}: מסנן אור כחול, נוח יותר לצפייה בטלפון בלילה.
  \item \textbf{מצב כהה}: חוסך בסוללה (יעיל במיוחד למסכי OLED), ואינו מסנוור בלילה. אפליקציות רבות תומכות גם הן במצב כהה.
  \item \textbf{מצב כהה}: ניתן לכוונון בחלק ממכשירי האנדרואיד, ככל שהוא גבוה יותר (למשל 120Hz) הגלילה חלקה יותר.
  \item \textbf{נעילה אוטומטית}: הגדר ל-15 או 30 שניות, בטוח וחסכוני בסוללה.
  \item \textbf{נעילה אוטומטית}: האותיות קטנות מדי? הגדל קצת.
  \item \textbf{גודל גופן}החלף לטפט יפה (סטטי או דינמי), ומצב הרוח משתנה מיד. לטלפוני אנדרואיד מותגיים יש לרוב חנות ערכות נושא, המאפשרת להחליף חבילה שלמה של אייקונים וסגנון.
\end{itemize}

\booksection{טפטים, ערכות נושא ומסך הבית3. טפטים, ערכות נושא ומסך הבית}

: מזג אוויר, יומן, בקרת מוזיקה על מסך הבית – גישה ישירה בצעד אחד. נתמך גם ב-iOS וגם באנדרואיד, לחץ לחיצה ארוכה על אזור ריק במסך הבית להוספה.

\textbf{4. הגדרות יעילות: הפוך את הטלפון ל“מוח החיצוני” שלך}אלה תכונות שברגע שמגדירים אותן “אי אפשר לחזור אחורה”, כל אחת מהן יכולה לשפר משמעותית את היעילות היומיומית שלך.

\booksection{הגדרות יעילות: הפוך את הטלפון ל“מוח החיצוני” שלך4. הגדרות יעילות: הפוך את הטלפון ל“מוח החיצוני” שלך}

רעיון מרכזי

\booksubsection{מצב מיקוד / 'נא לא להפריע': החזר את תשומת הלב שנגנבה על ידי ההתראות}

\begin{itemize}
  \item \textbf{המלצות שימוש}:
  \item \textbf{· מצב עבודה}: אפשר רק לאנשי קשר "מסומנים בכוכב" ולאפליקציות חיוניות;\newline ·\textbf{מצב שינה}: מופעל אוטומטית, מאפשר רק את השעון המעורר;\newline · כושר / קריאה / פגישות: התאמה אישית לפי הפעילות הנוכחית.\newline שימוש מתקדם\newline : השתמש באוטומציה + מיקום / זמן כדי להפעיל מצבים אוטומטית. לדוגמה: קיצורי דרך ב-iOS יכולים להגדיר "הגעה למשרד מפעילה אוטומטית מצב עבודה".
  \item \textbf{2. ניהול התראות: תן לנקודות האדומות להפסיק להיות "שיחות טלפון קטלניות"}כבה התראות שיווקיות
\end{itemize}

\booksubsection{ניהול התראות: תן לנקודות האדומות להפסיק להיות “שיחות טלפון קטלניות”}

\begin{itemize}
  \item \textbf{מיין לפי חשיבות}: הצב תקשורת בין-אישית (שיחות, SMS, מסרים מידיים) בראש; חדשות, קניות, רשתות חברתיות ניתן להגדיר כ"התראות שקטות".
  \item \textbf{קיבוץ התראות}: קבץ התראות דומות יחד, כדי למנוע הצפה בעשרות הודעות בפתיחה אחת.
  \item \textbf{3. מחוות מהירות: הפוך את הפעולות למהירות בצעד אחד}צילום מסך
\end{itemize}

\booksubsection{מחוות מהירות: הפוך את הפעולות למהירות בצעד אחד}

\begin{itemize}
  \item \textbf{צילום מסך ארוך / צילום מסך בגלילה}: לאחר צילום מסך רגיל, לחץ על אפשרות "צילום מסך ארוך".
  \item \textbf{סריקת קוד בלחיצה אחת}: הוסף את סמל "סרוק" במרכז הבקרה / הגדרות מהירות.
  \item \textbf{הקשה על הגב (iPhone)}: ב-[הגדרות] > [נגישות] > [מגע] > [הקשה על הגב] קשור פעולות כמו "צילום מסך" או "נעילת מסך".
  \item \textbf{התעוררות קולית}: Siri / XiaoAi / XiaoYi / XiaoBu, פשוט תקרא להם כדי לבדוק מזג אוויר או לכוון שעון מעורר.
  \item \textbf{4. מקלדת: כלי היעילות הכי לא מוערך}בחר את המקלדת הנכונה
\end{itemize}

\booksubsection{מקלדת: כלי היעילות הכי לא מוערך}

\begin{itemize}
  \item \textbf{· Sogou / Baidu}: מאגר מילים גדול, עדכון מהיר של מילות סלנג;\newline ·\textbf{iFlytek}: שיעור הזיהוי הגבוה ביותר להקלדה קולית, כלי עזר לרישום פגישות;\newline · Gboard (גוגל): פשוט, סנכרון חזק בין פלטפורמות, תומך בהקלדה בהחלקה;\newline · מקלדת מובנית iOS: הפרטיות הטובה ביותר, אינטגרציה עמוקה עם המערכת.\newline מתגים הכרחיים\newline : תיקון חכם, השלמת מילים, סנכרון מילון ענן, החלקה למעלה להזנת מספרים וסמלים.
  \item \textbf{טיפ לפרטיות}: היזהר בשימוש ב"השלמת משפטים" או "המלצות מותאמות אישית", העלולות להעלות את תוכן ההקלדה.
  \item \textbf{5. מסך בית וריבוי משימות: תן לזרימת העבודה לזרום}פיצול מסך / תמונה בתמונה
\end{itemize}

\booksubsection{מסך בית וריבוי משימות: תן לזרימת העבודה לזרום}

\begin{itemize}
  \item \textbf{חלון קטן / חלון צף}: WeChat / Douyin תומכים ב"מצב חלון קטן", השב וסגור מיד.
  \item \textbf{הצמד אפליקציות נפוצות}: שים רק את 4 האפליקציות בתדירות הגבוהה ביותר בשורת ה-Dock (בתחתית המסך), צמצם זמן חיפוש.
  \item \textbf{רשימת בדיקה ל"יום הראשון בעבודה"}לאחר קבלת טלפון חדש או עדכון גרסה גדול, מומלץ להקדיש 10 דקות לעבור על ההגדרות האלה:
\end{itemize}

\begin{tcbinfo}{(1) קוד נעילה + זיהוי ביומטרי (2) קיבוץ התראות + כיבוי התראות שיווק (3) גודל גופן / קצב רענון / מצב כהה}
(4) מקלדת ברירת מחדל / דפדפן / דוא“ל (5) אנשי קשר לשעת חירום / כרטיס רפואי (6) הגדרות גיבוי\newline 5. נגישות ופונקציות עזר: כוחה של טכנולוגיה למען הטוב\newline אל תיבהל מהשם ”פונקציות עזר“, התכונות האלה שימושיות לכולם. הן תוכננו במקור עבור אנשים עם מוגבלויות, אך הן שימושיות גם לאנשים רגילים – פציעה זמנית ביד, תינוק על הידיים, אוזניות מקולקלות ברכבת התחתית, משקפי קריאה לאיתור אותיות קטנות... כולן יכולות להציל את המצב. בסופו של דבר, עיצוב נגיש משקף את התפיסה של

\end{tcbinfo}

\booksection{נגישות ופונקציות עזר: כוחה של טכנולוגיה למען הטוב5. נגישות ופונקציות עזר: כוחה של טכנולוגיה למען הטוב}

: לאפשר לכל אחד ליהנות מחיי דיגיטל באופן שוויוני.\textbf{1. ה“עיניים” של משתמשים עם לקות ראייה – מדריך מלא לקורא מסך}קורא המסך הוא הכלי המרכזי למשתמשים עם לקות ראייה. הוא

\booksubsection{ה“עיניים” של משתמשים עם לקות ראייה – מדריך מלא לקורא מסך}

, והמשתמשים מפעילים את הטלפון באמצעות מחוות וקול. להלן שיטות השימוש ב-“VoiceOver” של iOS וב-TalkBack של אנדרואיד.\textbf{iOS VoiceOver}נתיב הפעלה

\begin{tcbinfo}{: [הגדרות] > [נגישות] > [VoiceOver], הפעל את המתג. ניתן גם ללחוץ שלוש פעמים על כפתור הצד (או כפתור הבית) להפעלה/כיבוי מהירים.}
\textbf{מחוות בסיסיות}:

החלקה באצבע אחת ימינה/שמאלה: מעבר בין רכיבים ממוקדים;

\begin{enumerate}
  \item הקשה כפולה באצבע אחת: הפעלת הרכיב הממוקד (שווה ערך ל"לחיצה");
  \item החלקה בשלוש אצבעות: מעבר בין דפים;
  \item הקשה משולשת באצבע אחת: הגדלה או בחירה (שונה מהקשה משולשת רגילה).
  \item הגדרות אישיות
  \item : [VoiceOver] > [דיבור], ניתן לכוון
\end{enumerate}

\textbf{קצב דיבור}(ברירת מחדל כ-50\%, ניתן להעלות ל-70-80\% לשמיעה חלקה יותר),\textbf{גובה צליל}, החלפת חבילות קול. iOS כוללת “דיבור משופר” ו“בקר קולי”, הראשון מפחית קולות נשימה והפוגות לשמיעה טבעית יותר.\textbf{מנוע דיבור}נתיב הפעלה

\end{tcbinfo}

\begin{tcbinfo}{: [הגדרות] > [נגישות] > [TalkBack], הפעל את המתג. או לחץ לחיצה ארוכה על כפתורי עוצמת השמע למשך 3 שניות להפעלה מהירה (נתמך בחלק מהמותגים).}
\textbf{מחוות בסיסיות}:

החלקה באצבע אחת ימינה/שמאלה: מעבר בין רכיבים ממוקדים;

\begin{enumerate}
  \item החלקה בשתי אצבעות למעלה/למטה/שמאלה/ימינה: גלילה;
  \item הקשה בשלוש אצבעות: כניסה למצב קריאה (הקראה רציפה);
  \item הקשה בארבע אצבעות: השהיה/המשך הקראה.
  \item הגדרות אישיות
  \item : [TalkBack] > [הגדרות טקסט לדיבור], ניתן לכוון
\end{enumerate}

\textbf{קצב דיבור},\textbf{גובה צליל}ו\textbf{מנוע דיבור}(מומלץ להתקין את מנוע הטקסט לדיבור של Google להקראה טבעית יותר).\textbf{2. הפוך אותיות קטנות לגדולות, הפוך מטושטש לברור – התאמות חזותיות}מחוות הגדלה

\end{tcbinfo}

\booksubsection{הפוך אותיות קטנות לגדולות, הפוך מטושטש לברור – התאמות חזותיות}

\booksubsection{מחוות הגדלה}

\begin{itemize}
  \item \textbf{פונקציות נפוצות:}להגדלה. ברירת מחדל זום מסך מלא, ניתן לבחור זום חלון (חלון הגדלה מקומי).\textbf{צביטה בשתי אצבעות מכווננת את ההגדלה (1.5x עד 15x)}, הקשה כפולה בשלוש אצבעות נועלת את יחס ההגדלה הנוכחי.\textbf{אנדרואיד}: [הגדרות] > [נגישות] > [הגדלה], הקש בשלוש אצבעות להגדלה, הקש שוב בשלוש אצבעות ליציאה; או
  \item \textbf{HarmonyOS}להפעלת הגדלה זמנית (עד פי 8).\textbf{טקסט בניגודיות גבוהה / טקסט מודגש}iOS
\end{itemize}

\booksubsection{טקסט בניגודיות גבוהה / טקסט מודגש}

\begin{itemize}
  \item \textbf{פונקציות נפוצות:}: [הגדרות] > [תצוגה] > [גודל גופן], ניתן לכוון לארבע רמות: "קטן, רגיל, גדול, גדול במיוחד". Huawei/Xiaomi מציעות בנוסף מתגי "מודגש" ו"טקסט בניגודיות גבוהה".
  \item \textbf{HarmonyOS}iOS
\end{itemize}

\booksubsection{מסנני צבע (ידידותי לעיוורי צבעים)}

\begin{itemize}
  \item \textbf{פונקציות נפוצות:}, וכן אפשרויות "גווני אפור" ו"היפוך צבעים". מומלץ לבצע בדיקת ראיית צבעים (מובנית במערכת) לפני בחירת המסנן המתאים.\textbf{אנדרואיד}: [הגדרות] > [נגישות] > [תיקון צבע], כולל
  \item \textbf{HarmonyOS}+ מתג "היפוך צבעים".\textbf{3. תיאור תמונות באמצעות AI: תן ל-AI להיות ה"עיניים" שלך}כשקורא מסך נתקל בתמונה, הגישה המסורתית היא לדלג או רק להקריא "תמונה". אבל עכשיו AI יכול
\end{itemize}

\booksubsection{תיאור תמונות באמצעות AI: תן ל-AI להיות ה“עיניים” שלך}

, דבר שימושי מאוד לזיהוי קוד אבטחה, קריאת תמונות וזיהוי תפריטים.\textbf{תמיכת תיאור תמונות AI לפי מערכת}iOS 15+

\begin{tcbinfo}{: ב-[הגדרות] > [נגישות] > [VoiceOver] > [תיאור קולי] הפעל "תיאור תמונות". המערכת תשתמש}
\begin{itemize}
  \item \textbf{במודל AI מקומי}ליצירת תיאור טקסט אוטומטי לתמונות (דורש iPhone 12 ומעלה).\textbf{אפליקציית TalkBack חדשה}: Google TalkBack בגרסאות חדשות יותר משלב
  \item \textbf{שאלות ותשובות על תמונות}, ניתן לשאול "מה יש בתמונה הזו" וה-AI עונה בזמן אמת (לא תלוי לחלוטין בגרסת מערכת אנדרואיד, עדכון אפליקציית TalkBack מספיק).\textbf{אפליקציות צד שלישי}: "Seeing AI" (iOS), "Lookout" (אנדרואיד), "XiaoYi BangBang" (Huawei) הן כולן עוזרות ראייה AI מצוינות, יכולות לזהות חפצים, לקרוא מסמכים ולזהות הבעות פנים.
  \item \textbf{תרחישי שימוש}נסה לשאול על תמונת רחוב "כמה אנשים יש בתמונה הזו? מה הם עושים?", או פתח אפליקציית משלוחי אוכל ותן ל-AI להקריא את שמות המנות מתמונת התפריט –
\end{itemize}

\end{tcbinfo}

\begin{tcbpractice}{תרחישי שימוש}
נסו לשאול על תמונת רחוב “כמה אנשים יש בתמונה הזו? מה הם עושים?”, או פתחו אפליקציית משלוחי אוכל ותנו ל-AI להקריא את שמות המנות מתמונת התפריט —\textbf{זו באמת “טכנולוגיה שהופכת את החיים לשוויוניים יותר”}.

\end{tcbpractice}

\booksubsection{ה“אוזניים” של משתמשים עם לקות שמיעה – מדריך מלא לכתוביות בזמן אמת}

, שימושית במיוחד למשתמשים עם לקות שמיעה בסביבות רועשות, בספריות ובתרחישי שפה זרה.\textbf{הפעלה ואופטימיזציה של כתוביות בזמן אמת}iOS

\begin{tcbinfo}{: [הגדרות] > [נגישות] > [כתוביות בזמן אמת], הפעל את המתג. המערכת תזהה את כל האודיו במכשיר (שיחות, וידאו, פודקאסטים, FaceTime, תזכירים קוליים) ותציג אוטומטית כתוביות בחלק העליון של המסך.}
\textbf{פונקציות נפוצות:}: לחץ על

\textbf{כפתור עוצמת השמע}הקש על סמל “כתוביות בזמן אמת” בחלון הקופץ; או הפעל ב-[הגדרות] > [נגישות] > [כתוביות בזמן אמת].\textbf{טיפים לשיפור דיוק}:

השתמשו ככל האפשר בסביבה שקטה, הרחק ממזגנים ומרעש מכונות;

\begin{enumerate}
  \item הדובר צריך לבטא בבירור ובקצב בינוני לדיוק מירבי;
  \item שיחות מרובות משתתפים מפחיתות דיוק — ניתן להשתמש במצב "הצג רק דובר אחד";
  \item תומך במנדרינית, קנטונזית, אנגלית ושפות נפוצות נוספות; בשפות פחות נפוצות הדיוק נמוך יותר;
  \item מוזיקת רקע או הפרעות קול אנושי במהלך הדיבור יגרמו לזיהוי מקוטע.
\end{enumerate}

\end{tcbinfo}

\booksubsection{ערוץ כפול בתרחישי שיחות: דיבור לטקסט + טקסט לדיבור}

(כמו אובדן קול פתאומי, התאוששות מניתוח גרון, גמגום חמור), הטלפון יכול לשמש כ“מתורגמן”.\textbf{תרגום שיחות בזמן אמת (TTY / RTT)}iOS

\begin{tcbinfo}{: [הגדרות] > [נגישות] > [RTT/TTY], לאחר הפעלה יוצג מקלדת על המסך במהלך שיחה, הדיבור של הצד השני יומר לטקסט בזמן אמת,}
\textbf{פונקציות נפוצות:}לצד השני.\textbf{אנדרואיד}: [הגדרות] > [נגישות] > [תמלול בזמן אמת] + [התראות קוליות]. “תמלול בזמן אמת” מציג את דיבור הצד השני כטקסט בזמן אמת; “התראות קוליות” ממירות את הקלט שלך להשמעה קולית כשדיבור אינו נוח.

\textbf{HarmonyOS}:

\textbf{תאימות}:

\begin{itemize}
  \item \textbf{TTY (טלפרינטר)}: דורש מכשיר TTY ייעודי, מוחלף בהדרגה על ידי RTT;
  \item \textbf{פונקציית}המרת דיבור לטקסט
  \item ב-WeChat/QQ היא גם חלופה: לחץ לחיצה ארוכה על הודעה קולית "המר לטקסט".\textbf{6. תכונות "קטנות אך שימושיות" לכולם}הקשה על הגב
\end{itemize}

\end{tcbinfo}

\booksubsection{תכונות “קטנות אך שימושיות” לכולם}

\begin{itemize}
  \item \textbf{הקשה כפולה/משולשת}להפעלת צילום מסך, נעילת מסך, מרכז בקרה ועוד. שימושי במיוחד כשהידיים עמוסות בחבילות.\textbf{מגע מסייע / כדור צף}: iOS [נגישות] > [מגע] > [מגע מסייע], אנדרואיד [נגישות] > [ניווט צף]. כפתור וירטואלי כתחליף למחוות או כפתורים פיזיים, מציל את המצב כשהכפתורים לא מגיבים או בתפעול ביד אחת.
  \item \textbf{בקרת מתג}: שליטה בטלפון באמצעות מתג חיצוני, הקשה בודדת, תנועות ראש,
  \item \textbf{מיועד למשתמשים עם מוגבלות פיזית}. iOS [נגישות] > [בקרת מתג], אנדרואיד [נגישות] > [בקרת מתג].\textbf{זיהוי קול}(iOS): [נגישות] > [זיהוי קול], לאחר הפעלה הטלפון "יקשיב" לפעמון דלת, בכי תינוק, אזעקת אש, סירנות ועוד,
  \item \textbf{עם זיהוי – רטט + התראה על המסך}, בעל ערך רב למשתמשים עם לקות שמיעה ובזמן שינה.\textbf{התראה באמצעות הבהוב LED}: iOS [נגישות] > [אודיו/ויזואלי] > [הבהוב LED להתראות], אנדרואיד [הגדרות] > [התראות] > [התראת פלאש]. הפלאש מהבהב בשיחות והתראות,
  \item \textbf{אודיו מונו}: [הגדרות] > [נגישות] > [אודיו] > [אודיו מונו] — ממזג את ערוצי השמע הימני והשמאלי. משתמשים עם ליקוי באוזן אחת לא יאבדו את חצי המידע השני.
  \item \textbf{פילוסופיית עיצוב הנגישות של אפל}אפל משקיעה משאבים עצומים בתחום הנגישות, ומעדכנת ומוסיפה פונקציות נגישות מדי שנה. לאתר הרשמי של אפל יש דף נגישות ייעודי המפרט את שיטות השימוש בכל תכונה. תכונות הנגישות של iOS בדרך כלל מלאות יותר מאלה של אנדרואיד עם חוויה חלקה יותר, וזו סיבה חשובה לכך שמשתמשים רבים עם מוגבלויות בוחרים באפל.
\end{itemize}

\begin{tcbinfo}{הקדש 30 שניות להתנסות}
מצא זמן פנוי ועשה את ההתנסויות הבאות:

\end{tcbinfo}

\begin{tcbpractice}{הפעל}
VoiceOver / TalkBack

\begin{enumerate}
  \item , נסה לדפדף ב-WeChat בהחלקה בשלוש אצבעות בלי להסתמך על הראייה;\textbf{הפעל}כתוביות בזמן אמת
  \item פתח\textbf{השתמש}בהקשה בשלוש אצבעות
  \item להגדלת טקסט ל-200\%, חווה את חיי היומיום של משתמשים עם לקות ראייה;\textbf{כבה את כל הפונקציות הקודמות, וחשוב "אם הייתי צריך את התכונות האלה לתמיד, איך הייתי רוצה שהן יעוצבו".}פרק 6: הגן על חייך הדיגיטליים – אבטחה ופרטיות
  \item אבטחת טלפון = אבטחת פרטיות + אבטחת רכוש. בעידן שבו נתונים הם נכס, להגן על הטלפון זה להגן על עצמך.
\end{enumerate}

\end{tcbpractice}
